% ============================================================================
% VALUE ARCHITECTURE: DATA AND EXECUTION IN A UNIFIED FRAME
% ============================================================================
\subsection{Unified Data--Execution Architecture}
\label{sec:value_architecture}

The central design decision of the system is the elimination of the
boundary between stored data and executable code.  A single Value frame
serves simultaneously as a record in memory and as a self-contained
computation unit.  This subsection describes how the frame layout
supports both roles and how their interaction gives rise to computation.

\subsubsection{The Data Plane}

The data region (words 0--59) holds the raw content of a Value: packed
token identifiers derived from input bytes, together with structural
metadata---a unique frame identifier and bidirectional linking pointers
that allow frames to reference their neighbors in a sequence.  The
affinity bitmask (word~64) encodes a 64-bit topological signature that
can be used to direct a frame toward compatible Regions or processing
paths.  The gossip field (words 66--69) provides a 256-bit fingerprint
summarizing the frame's content for lightweight similarity checks during
distributed routing.

From the perspective of the data plane, a Value is a passive container:
a fixed-size record whose fields can be read, compared, and combined
with the fields of other Values through bitwise operations.

\subsubsection{The Execution Plane}

The same frame also contains a complete execution context.  Seven
general-purpose registers (words 71--77) and a program counter (word~78)
define the machine state, while the program region (words 79--127) holds
up to 98 packed instructions.  When a frame is submitted to the compute
backend, the backend reads the instruction at the current program
counter, decodes its four-bit opcode as one of the sixteen boolean truth
tables, applies the operation across the specified source and
destination regions, and advances the program counter.

Execution halts when the program counter reaches an instruction slot
containing the zero opcode (the constant-false function) beyond the
first slot, or when the program counter exceeds the program region.
The backend reports the exit condition---normal exhaustion, explicit
halt, or invalid program word---through a status field embedded in the
frame's state region (word~63).

\subsubsection{Interaction Between Planes}

The data and execution planes occupy the same address space within the
frame.  An instruction in the program region can target any word in the
data region, and conversely, a boolean operation applied to the program
region can rewrite an instruction.  This bidirectional access has three
consequences:

First, a frame can \emph{inspect its own data}.  A program can read
token identifiers, compare affinity bitmasks, or accumulate statistics
over the data region using boolean operations, then store the result in
a register for use by subsequent instructions.

Second, a frame can \emph{modify its own program}.  Because program
words are addressable data, an instruction can overwrite a later
instruction, implementing conditional behavior without explicit branch
primitives.  The only control-flow mechanism is mutation of the program
counter or the program region itself.

Third, when two frames interact through the compute backend's
\texttt{UniversalBitwise} operation, each frame can read from the
other's data and program regions.  This enables one frame's program to
operate on another frame's data, or to copy instructions from one
frame to another---the mechanism underlying the viral firmware, which
propagates its program into a partner frame.

The result is a computational model in which storage and execution are
not separate subsystems connected by a bus, but overlapping views of the
same binary structure.
